
\begin{register}{H}{mvendorid}{0xF11}\label{regdesc:MVENDORID}
\regfield{MVENDORID}{32}{0}{{0x00000612}}%
\reglabel{Reset}\regnewline%
\begin{regdesc}\begin{reglist}
\label{fielddesc:MVENDORID}\item [MVENDORID] 表示供应商编号。 (RO)
\end{reglist}\end{regdesc}
\end{register}

\begin{register}{H}{marchid}{0xF12}\label{regdesc:MARCHID}
\regfield{MARCHID}{32}{0}{{0x80000002}}%
\reglabel{Reset}\regnewline%
\begin{regdesc}\begin{reglist}
\label{fielddesc:MARCHID}\item [MARCHID] 表示架构编号。 (RO)
\end{reglist}\end{regdesc}
\end{register}

\begin{register}{H}{mimpid}{0xF13}\label{regdesc:MIMPID}
\regfield{MIMPID}{32}{0}{{0x00000002}}%
\reglabel{Reset}\regnewline%
\begin{regdesc}\begin{reglist}
\label{fielddesc:MIMPID}\item [MIMPID] 表示硬件实现编号。(RO)
\end{reglist}\end{regdesc}
\end{register}

\begin{register}{H}{mhartid}{0xF14}\label{regdesc:MHARTID}
\regfield{MHARTID}{32}{0}{{0x00000000}}%
\reglabel{Reset}\regnewline%
\begin{regdesc}\begin{reglist}
\label{fielddesc:MHARTID}\item [MHARTID] 表示硬件线程编号。(RO)
\end{reglist}\end{regdesc}
\end{register}

\begin{register}{H}{mstatus}{0x300}\label{regdesc:MSTATUS}
\regfield{(reserved)}{10}{22}{{0x000}}%
\regfield{TW}{1}{21}{{0}}%
\regfield{(reserved)}{8}{13}{{0x00}}%
\regfield{MPP}{2}{11}{{0x0}}%
\regfield{(reserved)}{3}{8}{{0x0}}%
\regfield{MPIE}{1}{7}{{0}}%
\regfield{(reserved)}{2}{5}{{0x0}}%
\regfield{UPIE}{1}{4}{{0}}%
\regfield{MIE}{1}{3}{{0}}%
\regfield{(reserved)}{2}{1}{{0x0}}%
\regfield{UIE}{1}{0}{{0}}%
\reglabel{Reset}\regnewline%
\begin{regdesc}\begin{reglist}
\label{fielddesc:MSTATUSUIE}\item [UIE] 写 1 使能全局用户模式中断。 (R/W)
\label{fielddesc:MSTATUSMIE}\item [MIE] 写 1 使能全局机器模式中断。(R/W)
\label{fielddesc:MSTATUSUPIE}\item [UPIE] 写 1 使能（异常）之前的用户模式中断。(R/W)
\label{fielddesc:MSTATUSMPIE}\item [MPIE] 写 1 使能（异常）之前的机器模式中断。 (R/W)
\label{fielddesc:MSTATUSMPP}\item [MPP] 配置（异常）之前的机器特权模式。\\
0x0：用户模式\\
0x3：机器模式\\
注意：仅低位可写。由于高位直接绑定低位，写入高位将被忽略。\\
(R/W)
\label{fielddesc:MSTATUSTW}\item [TW] 配置在用户模式下执行 WFI（等待中断）指令是否导致非法指令异常。\\
0：不导致非法指令异常\\
1：导致非法指令异常\\
(R/W)
\end{reglist}\end{regdesc}
\end{register}

\begin{register}{H}{misa}{0x301}\label{regdesc:MISA}
\regfield{MXL}{2}{30}{{0x1}}%
\regfield{(reserved)}{4}{26}{{0x0}}%
\regfield{Z}{1}{25}{{0}}%
\regfield{Y}{1}{24}{{0}}%
\regfield{X}{1}{23}{{0}}%
\regfield{W}{1}{22}{{0}}%
\regfield{V}{1}{21}{{0}}%
\regfield{U}{1}{20}{{1}}%
\regfield{T}{1}{19}{{0}}%
\regfield{S}{1}{18}{{0}}%
\regfield{R}{1}{17}{{0}}%
\regfield{Q}{1}{16}{{0}}%
\regfield{P}{1}{15}{{0}}%
\regfield{O}{1}{14}{{0}}%
\regfield{N}{1}{13}{{0}}%
\regfield{M}{1}{12}{{1}}%
\regfield{L}{1}{11}{{0}}%
\regfield{K}{1}{10}{{0}}%
\regfield{J}{1}{9}{{0}}%
\regfield{I}{1}{8}{{1}}%
\regfield{H}{1}{7}{{0}}%
\regfield{G}{1}{6}{{0}}%
\regfield{F}{1}{5}{{0}}%
\regfield{E}{1}{4}{{0}}%
\regfield{D}{1}{3}{{0}}%
\regfield{C}{1}{2}{{1}}%
\regfield{B}{1}{1}{{0}}%
\regfield{A}{1}{0}{{1}}%
\reglabel{Reset}\regnewline%
\begin{regdesc}\begin{reglist}
\label{fielddesc:MISAMXL}\item [MXL] 机器 XLEN = 1（32 位）。(RO)
\label{fielddesc:MISAZ}\item [Z] 保留 = 0。(RO)
\label{fielddesc:MISAY}\item [Y] 保留 = 0。(RO)
\label{fielddesc:MISAX}\item [X] 非标准扩展 = 0。(RO)
\label{fielddesc:MISAW}\item [W] 保留 = 0。(RO)
\label{fielddesc:MISAV}\item [V] 保留 = 0。(RO)
\label{fielddesc:MISAU}\item [U] 实现用户模式 = 1。(RO)
\label{fielddesc:MISAT}\item [T] 保留 = 0。(RO)
\label{fielddesc:MISAS}\item [S] 实现监督模式 = 0。(RO)
\label{fielddesc:MISAR}\item [R] 保留 = 0。(RO)
\label{fielddesc:MISAQ}\item [Q] 四精度浮点扩展 = 0。(RO)
\label{fielddesc:MISAP}\item [P] 保留 = 0。(RO)
\label{fielddesc:MISAO}\item [O] 保留 = 0。(RO)
\label{fielddesc:MISAN}\item [N] 支持用户级别中断 = 0。(RO)
\label{fielddesc:MISAM}\item [M] 整数乘除法标准扩展 = 1。(RO)
\label{fielddesc:MISAL}\item [L] 保留 = 0。(RO)
\label{fielddesc:MISAK}\item [K] 保留 = 0。(RO)
\label{fielddesc:MISAJ}\item [J] 保留 = 0。(RO)
\label{fielddesc:MISAI}\item [I] RV32I 基本 ISA = 1。(RO)
\label{fielddesc:MISAH}\item [H] 虚拟机管理程序扩展 = 0。(RO)
\label{fielddesc:MISAG}\item [G] 其他标准扩展 = 0。(RO)
\label{fielddesc:MISAF}\item [F] 单精度浮点扩展 = 0。(RO)
\label{fielddesc:MISAE}\item [E] RV32E 基本 ISA = 0。(RO)
\label{fielddesc:MISAD}\item [D] 双精度浮点扩展 = 0。(RO)
\label{fielddesc:MISAC}\item [C] 压缩标准扩展 = 1。(RO)
\label{fielddesc:MISAB}\item [B] 保留 = 0。(RO)
\label{fielddesc:MISAA}\item [A] 原子标准扩展 = 1。(RO)
\end{reglist}\end{regdesc}
\end{register}

\begin{register}{H}{mideleg}{0x303}\label{regdesc:MIDELEG}
\regfield{MIDELEG}{32}{0}{{0x00000111}}%
\reglabel{Reset}\regnewline%
\begin{regdesc}\begin{reglist}
\label{fielddesc:MIDELEG}\item [MIDELEG] 配置是否委托相应中断给用户模式。以下中断默认委托给 U 模式：\\
Bit 0：用户软件中断 (CLINT)\\
Bit 4：用户定时器中断 (CLINT)\\
Bit 8：用户外部中断\\
可根据运行时的实际需求修改默认委托给用户模式的中断。\\
(R/W)
\end{reglist}\end{regdesc}
\end{register}

\begin{register}{H}{mie}{0x304}\label{regdesc:MIE}
\regfield{MXIE[31:8]}{24}{8}{{0x0}}%
\regfield{MTIE}{1}{7}{{0x0}}%
\regfield{MXIE[6:5]}{2}{5}{{0x0}}%
\regfield{UTIE}{1}{4}{{0x0}}%
\regfield{MSIE}{1}{3}{{0x0}}%
\regfield{MXIE[2:1]}{2}{1}{{0x0}}%
\regfield{USIE}{1}{0}{{0x0}}%
\reglabel{Reset}\regnewline%
\begin{regdesc}\begin{reglist}
\label{fielddesc:MIEUSIE}\item [USIE] 写 1 使能用户软件中断。(R/W)
\label{fielddesc:MIEMSIE}\item [MSIE] 写 1 使能机器软件中断。(R/W)
\label{fielddesc:MIEUTIE}\item [UTIE] 写 1 使能用户定时器中断。 (R/W)
\label{fielddesc:MIEMTIE}\item [MTIE] 写 1 使能机器定时器中断。(R/W)
\label{fielddesc:MIEMXIE}\item [MXIE] 写 1 使能 28 个外部中断。(R/W)
\end{reglist}\end{regdesc}
\end{register}

\begin{register}{H}{mtvec}{0x305}\label{regdesc:MTVEC}
\regfield{BASE}{24}{8}{{0x000000}}%
\regfield{(reserved)}{6}{2}{{0x00}}%
\regfield{MODE}{2}{0}{{0x1}}%
\reglabel{Reset}\regnewline%
\begin{regdesc}\begin{reglist}
\label{fielddesc:MTVECMODE}\item [MODE] 表示机器模式中断是否为向量模式。仅支持向量模式 \textbf{0x1}。(RO)
\label{fielddesc:MTVECBASE}\item [BASE] 配置异常向量基址的高 24 位，地址为 256 字节对齐。 (R/W)
\end{reglist}\end{regdesc}
\end{register}

\begin{register}{H}{mscratch}{0x340}\label{regdesc:MSCRATCH}
\regfield{MSCRATCH}{32}{0}{{0x00000000}}%
\reglabel{Reset}\regnewline%
\begin{regdesc}\begin{reglist}
\label{fielddesc:MSCRATCH}\item [MSCRATCH] 配置用户自定义的机器暂存信息。(R/W)
\end{reglist}\end{regdesc}
\end{register}

\begin{register}{H}{mepc}{0x341}\label{regdesc:MEPC}
\regfield{MEPC}{32}{0}{{0x00000000}}%
\reglabel{Reset}\regnewline%
\begin{regdesc}\begin{reglist}
\label{fielddesc:MEPC}\item [MEPC] 配置机器陷阱/异常程序计数器。当 CPU 遇到异常时，此域将自动更新为 CPU 将要执行的指令的地址。 (R/W)
\end{reglist}\end{regdesc}
\end{register}

\begin{register}{H}{mcause}{0x342}\label{regdesc:MCAUSE}
\regfield{Interrupt Flag}{1}{31}{{0}}%
\regfield{(reserved)}{26}{5}{{0x0000000}}%
\regfield{Exception Code}{5}{0}{{0x00}}%
\reglabel{Reset}\regnewline%
\begin{regdesc}\begin{reglist}
\label{fielddesc:MCAUSECODE}\item [Exception Code] CPU 进入异常时，此域将自动更新为最近的异常或中断的唯一 ID。可能的异常 ID：\\
0x1：PMP 指令访问错误\\
0x2：非法指令\\
0x3：硬件断点/观察点或 EBREAK\\
0x5：PMP 读存储器访问错误\\
0x6：写存储器地址或 AMO 地址未对齐\\
0x7：PMP 写存储器访问错误\\
0x8：用户模式环境调用 (ECALL)\\
0xb：机器模式环境调用\\
其他值：保留\\
注意：异常 ID 0x0（指令地址非对齐）不存在，因为 CPU 在取指时始终屏蔽地址的最低位。\\
(R/W)
\label{fielddesc:MCAUSEINT}\item [Interrupt Flag]  CPU 进入异常时，此标志位将自动更新。\\
如果被置位，则表示最近的陷阱是由中断引起。在异常情况下保持为 0。\\
注意：中断控制器将中断编号 1-2、5-6、8-31 全部用于外部中断源，而 RISC-V 标准则为内核的内部中断源预留了编号 0-15。本地中断源 (CLINT) 使用 0、3、4、7 等保留编号。\\
(R/W) 
\end{reglist}\end{regdesc}
\end{register}

\begin{register}{H}{mtval}{0x343}\label{regdesc:MTVAL}
\regfield{MTVAL}{32}{0}{{0x00000000}}%
\reglabel{Reset}\regnewline%
\begin{regdesc}\begin{reglist}
\label{fielddesc:MTVAL}\item [MTVAL] 配置机器陷阱值。将自动更新为与异常有关的数据，该数据可能有助于处理该异常。根据异常编号有以下解读：\\
0x1：指令虚拟地址错误\\
0x2：指令 opcode 错误\\
0x5：存储器读操作的数据地址错误\\
0x7：存储器写操作的数据地址错误\\
注意：该寄存器不支持其他异常 ID 和中断。\\
(R/W)
\end{reglist}\end{regdesc}
\end{register}

\begin{register}{H}{mip}{0x344}\label{regdesc:MIP}
\regfield{MXIP[31:8]}{24}{8}{{0x0}}%
\regfield{MTIP}{1}{7}{{0x0}}%
\regfield{MXIP[6:5]}{2}{5}{{0x0}}%
\regfield{UTIP}{1}{4}{{0x0}}%
\regfield{MSIP}{1}{3}{{0x0}}%
\regfield{MXIP[2:1]}{2}{1}{{0x0}}%
\regfield{USIP}{1}{0}{{0x0}}%
\reglabel{Reset}\regnewline%
\begin{regdesc}\begin{reglist}
\label{fielddesc:MIPUSIP}\item [USIP] 配置用户软件中断的等待状态。\\
0：不等待\\
1：等待\\
(R/W)
\label{fielddesc:MIPMSIP}\item [MSIP] 配置机器软件中断的等待状态。 \\
0：不等待\\
1：等待\\
(R/W)
\label{fielddesc:MIPUTIP}\item [UTIP] 配置用户定时器中断的等待状态。 \\
0：不等待\\
1：等待\\
(R/W)
\label{fielddesc:MIPMTIP}\item [MTIP] 配置机器定时器中断的等待状态。 \\
0：不等待\\
1：等待\\
(R/W)
\label{fielddesc:MIPMXIP}\item [MXIP] 配置 28 个外部中断的等待状态。 \\
0：不等待\\
1：等待\\
(R/W)
\end{reglist}\end{regdesc}
\end{register}

\begin{register}{H}{ustatus}{0x300}\label{regdesc:USTATUS}
\regfield{(reserved)}{27}{5}{{0x0000000}}%
\regfield{UPIE}{1}{4}{{0}}%
\regfield{(reserved)}{3}{1}{{0x0}}%
\regfield{UIE}{1}{0}{{0}}%
\reglabel{Reset}\regnewline%
\begin{regdesc}\begin{reglist}
\label{fielddesc:USTATUSUIE}\item [UIE] 写 1 使能全局用户模式中断。 (R/W)
\label{fielddesc:USTATUSUPIE}\item [UPIE] 写 1 使能（异常）之前的用户中断。(R/W)
\end{reglist}\end{regdesc}
\end{register}

\begin{register}{H}{uie}{0x004}\label{regdesc:UIE}
\regfield{UXIE[31:8]}{24}{8}{{0x0}}%
\regfield{(reserved)}{1}{7}{{0x0}}%
\regfield{UXIE[6:5]}{2}{5}{{0x0}}%
\regfield{UTIE}{1}{4}{{0x0}}%
\regfield{(reserved)}{1}{3}{{0x0}}%
\regfield{UXIE[2:1]}{2}{1}{{0x0}}%
\regfield{USIE}{1}{0}{{0x0}}%
\reglabel{Reset}\regnewline%
\begin{regdesc}\begin{reglist}
\label{fielddesc:UIEUSIE}\item [USIE] 写 1 使能用户软件中断。(R/W)
\label{fielddesc:UIEUTIE}\item [UTIE] 写 1 使能用户定时器中断。(R/W)
\label{fielddesc:UIEUXIE}\item [UXIE] 写 1 使能 28 个委托给用户模式的外部中断。 (R/W)
\end{reglist}\end{regdesc}
\end{register}

\begin{register}{H}{utvec}{0x005}\label{regdesc:UTVEC}
\regfield{BASE}{24}{8}{{0x000000}}%
\regfield{(reserved)}{6}{2}{{0x00}}%
\regfield{MODE}{2}{0}{{0x1}}%
\reglabel{Reset}\regnewline%
\begin{regdesc}\begin{reglist}
\label{fielddesc:UTVECMODE}\item [MODE] 表示用户模式中断是否为向量模式，仅支持向量模式 \textbf{0x1}。 (RO)
\label{fielddesc:UTVECBASE}\item [BASE] 配置异常向量基址的高 24 位，地址为 256 字节对齐。 (R/W)
\end{reglist}\end{regdesc}
\end{register}

\begin{register}{H}{uscratch}{0x040}\label{regdesc:USCRATCH}
\regfield{USCRATCH}{32}{0}{{0x00000000}}%
\reglabel{Reset}\regnewline%
\begin{regdesc}\begin{reglist}
\label{fielddesc:USCRATCH}\item [USCRATCH] 配置用户自定义的机器暂存信息。 (R/W)
\end{reglist}\end{regdesc}
\end{register}

\begin{register}{H}{uepc}{0x041}\label{regdesc:UEPC}
\regfield{UEPC}{32}{0}{{0x00000000}}%
\reglabel{Reset}\regnewline%
\begin{regdesc}\begin{reglist}
\label{fielddesc:UEPC}\item [UEPC] 配置用户异常程序计数器。当 CPU 遇到异常时，此域将自动更新为 CPU 将要执行的指令的地址。(R/W)
\end{reglist}\end{regdesc}
\end{register}

\begin{register}{H}{ucause}{0x042}\label{regdesc:UCAUSE}
\regfield{Interrupt Flag}{1}{31}{{0}}%
\regfield{(reserved)}{26}{5}{{0x0000000}}%
\regfield{Exception Code}{5}{0}{{0x00}}%
\reglabel{Reset}\regnewline%
\begin{regdesc}\begin{reglist}
\label{fielddesc:UCAUSECODE}\item [Interrupt ID] 进入异常时，此域将自动更新为最近的异常或中断的唯一 ID。(R/W)
\label{fielddesc:UCAUSEINT}\item [Interrupt Flag] 此标志将始终设置，因为 CPU 只能由于用户模式中断而进入陷阱，因为不支持异常委托。  (R/W)
\end{reglist}\end{regdesc}
\end{register}

\begin{register}{H}{uip}{0x044}\label{regdesc:UIP}
\regfield{UXIP[31:8]}{24}{8}{{0x0}}%
\regfield{(reserved)}{1}{7}{{0x0}}%
\regfield{UXIP[5:4]}{2}{5}{{0x0}}%
\regfield{UTIP}{1}{4}{{0x0}}%
\regfield{(reserved)}{1}{3}{{0x0}}%
\regfield{UXIP[2:1]}{2}{1}{{0x0}}%
\regfield{USIP}{1}{0}{{0x0}}%
\reglabel{Reset}\regnewline%
\begin{regdesc}\begin{reglist}
\label{fielddesc:UIPUSIP}\item [USIP] 配置用户软件中断的等待状态。 \\
0：不等待\\
1：等待\\
(R/W)
\label{fielddesc:UIPUTIP}\item [UTIP] 配置用户定时器中断的等待状态。\\
0：不等待\\
1：等待\\
(R/W)
\label{fielddesc:UIPUXIP}\item [UXIP] 配置 28 个委托给用户模式的外部中断。\\
0：不等待\\
1：等待\\
(R/W)
\end{reglist}\end{regdesc}
\end{register}

\begin{register}{H}{mpcer}{0x7E0}\label{regdesc:MPCER}
\regfield{(reserved)}{21}{11}{{0x000}}%
\regfield{INST\_COMP}{1}{10}{{0}}%
\regfield{(BRANCH\_TAKEN}{1}{9}{{0}}%
\regfield{BRANCH}{1}{8}{{0}}%
\regfield{JMP\_UNCOND}{1}{7}{{0}}%
\regfield{STORE}{1}{6}{{0}}%
\regfield{LOAD}{1}{5}{{0}}%
\regfield{IDLE}{1}{4}{{0}}%
\regfield{JMP\_HAZARD}{1}{3}{{0}}%
\regfield{LD\_HAZARD}{1}{2}{{0}}%
\regfield{INST}{1}{1}{{0}}%
\regfield{CYCLE}{1}{0}{{0}}%
\reglabel{Reset}\regnewline%
\begin{regdesc}\begin{reglist}
\label{fielddesc:INST_COMP}\item [INST\_COMP] 计数压缩指令。(R/W)
\label{fielddesc:BRANCH_TAKEN}\item [BRANCH\_TAKEN] 计数分支跳转。(R/W)  
\label{fielddesc:BRANCH}\item [BRANCH] 计数分支。 (R/W) 
\label{fielddesc:JUMP_UNCONDITIONAL}\item [JMP\_UNCOND] 计数无条件跳转。 (R/W) 
\label{fielddesc:STORE}\item [STORE] 计数存储器写操作。 (R/W) 
\label{fielddesc:LOAD}\item  [LOAD] 计数存储器读操作。 (R/W) 
\label{fielddesc:IDLE}\item  [IDLE] 计数 IDLE 周期。 (R/W) 
\label{fielddesc:JUMP_HAZARDS}\item  [JMP\_HAZARD] 计数跳转冒险。(R/W) 
\label{fielddesc:LOAD_HAZARDS}\item  [LD\_HAZARD] 计数存储器读操作冒险。(R/W) 
\label{fielddesc:INSTRUCTION}\item  [INST] 计数指令。(R/W) 
\label{fielddesc:CYCLES}\item  [CYCLE] 计数时钟周期，WFI 模式下周期计数不增加。\\
注意：每个位选择一个特定事件由计数器递增计数。如果多个事件被选择且同时发生，则计数器只递增 1。\\
(R/W)

\end{reglist}\end{regdesc}
\end{register}

\begin{register}{H}{mpcmr}{0x7E1}\label{regdesc:MPCMR}
\regfield{(reserved)}{30}{2}{{0}}%
\regfield{COUNT\_SAT}{1}{1}{{1}}%
\regfield{COUNT\_EN}{1}{0}{{1}}%
\reglabel{Reset}\regnewline%
\begin{regdesc}\begin{reglist}
\label{fielddesc:COUNTER_SATURATION}\item [COUNT\_SAT] 配置计数器饱\\
0：超出最大值上溢\\
1：超出最大值暂停\\
(R/W)
\label{fielddesc:COUNTER_ENABLE}\item [COUNT\_EN] 配置 是否使能计数器。\\
0：不使能\\
1：使能\\
(R/W)
\end{reglist}\end{regdesc}
\end{register}

\begin{register}{H}{mpccr}{0x7E2}\label{regdesc:MPCCR}
\regfield{MPCCR}{32}{0}{{0x00000000}}%
\reglabel{Reset}\regnewline%
\begin{regdesc}\begin{reglist}
\label{fielddesc:MPCCR}\item [MPCCR] 表示机器性能计数器数值。(R/W)
\end{reglist}\end{regdesc}
\end{register}